\section{Algoritmos y técnicas de entrenamiento}

\subsection{Bibliotecas recomendadas de fine-tuning}

El presentador recomendó tres bibliotecas principales de fine-tuning:

\begin{table}[H]
\centering
\caption{Comparación de las principales bibliotecas de fine-tuning}
\begin{tabular}{llp{7cm}}
\toprule
\textbf{Herramienta} & \textbf{Desarrollador} & \textbf{Características} \\
\midrule
TRL & Hugging Face & Basada en Transformers, con la mayor variedad de algoritmos y actualizada con la investigación más reciente \\
Axolotl & Comunidad & Basada en TRL, impulsada por configuración YAML, fácil de compartir y reproducir configuraciones de entrenamiento \\
Unsloth & Comunidad & Optimizada para una sola GPU, con una eficiencia muy alta e incorpora funciones prácticas como la cuantización \\
\bottomrule
\end{tabular}
\end{table}

\subsection{Las tres técnicas de entrenamiento de SFT}

\begin{importantbox}{Full Fine-Tuning vs. LoRA vs. QLoRA}
\begin{itemize}
    \item \textbf{Full Fine-Tuning}: carga el modelo con precisión completa y entrena el 100\% de los parámetros. Ofrece la mayor calidad, pero requiere una cantidad enorme de VRAM y, por lo general, un clúster de GPU completo.
    \item \textbf{LoRA (Low-Rank Adaptation)}: congela los pesos originales del modelo y solo entrena las matrices adicionales de adaptador de bajo rango A y B. Solo necesita entrenar alrededor del 0.5\% de los parámetros, es rápido y requiere poca VRAM; \textbf{se recomienda como opción predeterminada}.
    \item \textbf{QLoRA (Quantized LoRA)}: primero cuantiza el modelo a una precisión de 4 bits y luego lo carga, para después aplicar LoRA. Reduce aún más los requisitos de VRAM, pero el entrenamiento es más lento y el rendimiento presenta una ligera caída (unos cuantos puntos porcentuales).
\end{itemize}
\end{importantbox}

Principio matemático de LoRA: para la matriz de pesos original $W \in \mathbb{R}^{d \times k}$, LoRA agrega una descomposición de bajo rango:

$$W' = W + \Delta W = W + BA$$

donde:
\begin{itemize}
    \item $B \in \mathbb{R}^{d \times r}$, $A \in \mathbb{R}^{r \times k}$
    \item $r \ll \min(d, k)$ es el rango (rank), y normalmente se elige entre 8 y 64
    \item Durante el entrenamiento solo se actualizan $A$ y $B$; el $W$ original permanece congelado
\end{itemize}

\begin{warningbox}{Full Fine-Tuning suele ser excesivo}
A menos que seas una empresa dedicada específicamente a entrenar LLM, LoRA es casi siempre la mejor opción. Full Fine-Tuning puede provocar sobreajuste, y la mejora en el rendimiento a menudo no justifica el costo computacional adicional. Solo si los requisitos de VRAM de LoRA siguen siendo demasiado altos conviene recurrir a QLoRA.
\end{warningbox}

\subsection{Algoritmos de Preference Alignment}

El campo de la alineación de preferencias cuenta con más de 100 algoritmos distintos, pero el presentador recomienda memorizar solo dos:

\subsubsection{PPO (Proximal Policy Optimization)}

PPO es el algoritmo clásico de RLHF y requiere cargar \textbf{tres modelos} simultáneamente:

\begin{enumerate}
    \item \textbf{Frozen Model} (política de referencia): una instantánea del modelo previa al entrenamiento, que se usa para calcular la KL divergence
    \item \textbf{Train Model} (política actual): el modelo que recibe el prompt y genera la respuesta
    \item \textbf{Reward Model}: califica la respuesta generada
\end{enumerate}

Flujo de entrenamiento: el Train Model genera una respuesta $\rightarrow$ el Reward Model la califica $\rightarrow$ se usa la KL divergence para restringir la magnitud de la actualización de la política $\rightarrow$ se actualizan los parámetros del Train Model.

La ventaja es que el límite superior de calidad es alto; la desventaja es que requiere tres modelos, es extremadamente costoso y difícil de ajustar.

\subsubsection{DPO (Direct Preference Optimization)}

DPO solo necesita \textbf{dos modelos} (si se usa LoRA, en realidad basta con cargar un modelo base más dos conjuntos de adaptadores):

$$\mathcal{L}_{\text{DPO}} = -\mathbb{E}\left[\log \sigma\left(\beta \log \frac{\pi_\theta(y_w|x)}{\pi_{\text{ref}}(y_w|x)} - \beta \log \frac{\pi_\theta(y_l|x)}{\pi_{\text{ref}}(y_l|x)}\right)\right]$$

donde:
\begin{itemize}
    \item $y_w$ es la chosen answer y $y_l$ es la rejected answer
    \item $\pi_\theta$ es la política en entrenamiento y $\pi_{\text{ref}}$ es la política de referencia
    \item $\beta$ controla la intensidad de la penalización por desviarse de la política de referencia
\end{itemize}

\begin{importantbox}{Recomendación práctica: DPO es mejor que PPO}
A menos que se cuente con recursos prácticamente ilimitados al nivel de OpenAI, se recomienda usar DPO. DPO es más rápido, más económico y más fácil de ajustar, con una calidad apenas ligeramente inferior a la de PPO. DPO no necesita un reward model independiente, sino que optimiza directamente sobre preference data, lo cual simplifica enormemente el flujo de entrenamiento.
\end{importantbox}

\subsection{Hiperparámetros clave del entrenamiento}

\begin{table}[H]
\centering
\caption{Hiperparámetros clave del Post-Training}
\begin{tabular}{lll}
\toprule
\textbf{Parámetro} & \textbf{Valor recomendado} & \textbf{Descripción} \\
\midrule
Learning Rate & $1 \times 10^{-5}$ $\sim$ $5 \times 10^{-5}$ & El parámetro más crítico; si es demasiado alto provoca un loss spike \\
Batch Size & Limitado por la VRAM & Se puede usar gradient accumulation para simular un batch grande \\
Max Length & Limitado por la VRAM & Determina la longitud máxima de la secuencia de entrada \\
Epochs & 3$\sim$5 & Número de veces que se recorre por completo el conjunto de entrenamiento \\
Optimizer & AdamW & Una base sólida que normalmente no hace falta cambiar \\
Attention & Flash Attention & Una implementación eficiente de la atención que acelera especialmente el procesamiento de secuencias largas \\
\bottomrule
\end{tabular}
\end{table}

\begin{knowledgebox}{Gradient Accumulation y Effective Batch Size}
Cuando la memoria de la GPU no basta para admitir el batch size deseado, se puede usar gradient accumulation: se acumulan los gradientes tras varios forward pass y luego se actualizan los pesos de una sola vez. De este modo, \textbf{effective batch size} = per-GPU batch size $\times$ accumulation steps $\times$ número de GPU, lo que permite simular un batch más grande sin aumentar la memoria.
\end{knowledgebox}

\subsection{Monitoreo del entrenamiento}

Durante el entrenamiento es necesario monitorear tres métricas fundamentales:

\begin{enumerate}
    \item \textbf{Training Loss}: debe disminuir de forma uniforme; si aparece un loss spike, por lo general significa que el learning rate es demasiado alto
    \item \textbf{Validation Loss}: se usa entre el 1 y el 2\% de los datos como conjunto de validación para monitorear si hay sobreajuste
    \item \textbf{Gradient Norm}: la magnitud del gradiente; demasiados spike pueden indicar un problema con los datos
\end{enumerate}

\begin{warningbox}{Diagnóstico de un Loss Spike}
Si al principio del entrenamiento el loss disminuye de forma estable, pero luego aparece de pronto un salto brusco (spike), la causa más probable es que el \textbf{learning rate sea demasiado alto}. La solución es reducir el learning rate, lo cual debería producir una curva de loss más uniforme. Ten cuidado de no confundir las pequeñas fluctuaciones del inicio del entrenamiento con un spike — las fluctuaciones iniciales de poca magnitud suelen ser normales.
\end{warningbox}

\subsection{Resumen de la sección}

Para SFT se recomienda usar LoRA como técnica de entrenamiento predeterminada, y para Preference Alignment se recomienda usar DPO. Entre los hiperparámetros clave, el learning rate es el más importante, y AdamW junto con Flash Attention es una opción predeterminada confiable. Durante el entrenamiento se debe monitorear de cerca el training loss, el validation loss y el gradient norm.

\newpage
